% p.579
$1 \times 1$ convolution $1 \times 1$ 卷积\quad 255\par
1-of-K coding 1-of-K 编码方案\quad 58, 119, 392\par
acceptance criterion 接受准则\quad 375, 378, 381\par
activation 激活\quad 15\par
activation function 激活函数\quad 15, 137, 156\par
active constraint 激活的约束\quad 528\par
AdaGrad AdaGrad\quad 193\par
Adam optimization Adam 算法\quad 193\par
adaptive rejection sampling 适应性拒绝采样\quad 370\par
adjacency matrix 邻接矩阵\quad 349\par
adjoint sensitivity method 伴随灵敏度\quad 473\par
adversarial attack 对抗攻击\quad 263\par
aggregation 聚合\quad 353\par
aleatoric uncertainty 偶然不确定性\quad 21\par
AlexNet AlexNet\quad 257\par
alpha family $\alpha$ 家族\quad 53\par
amortized inference 摊销推理\quad 487\par
ancestral sampling 祖先采样\quad 382\par
anchor 锚点\quad 166\par
annealed Langevin dynamics 退火朗之万动力学\quad 508\par
AR model, see autoregressive model AR 模型，参见“自回归模型”\par
area under the ROC curve ROC 曲线下面积\quad 130\par
artificial intelligence 人工智能\quad 1\par
attention 注意力\quad 306\par
attention head 注意力头\quad 313\par
audio data 音频数据\quad 339\par
auto-associative neural network, see autoencoder 自联想神经网络，参见“自编码器”\par
autoencoder 自编码器\quad 164, 479\par
automatic differentiation 自动微分\quad 19, 201, 211\par
autoregressive flow 自回归流\quad 470\par
autoregressive model 自回归模型\quad 5, 299, 323\par
average pooling 平均池化\quad 255\par
backpropagation 反向传播\quad 17, 201\par
backpropagation through time 通过时间反向传播\quad 325\par
bag of words 词袋\quad 322\par
bagging 装袋\quad 240\par
base distribution 基础分布\quad 465\par
basis function 基函数\quad 98, 138, 150\par
batch gradient descent 批量梯度下降\quad 185\par
batch learning 批量学习\quad 102\par
batch normalization 批量归一化\quad 196\par
Bayes net 贝叶斯网\quad 280\par

% p.580
Bayes’ theorem 贝叶斯定理\quad 26\par
Bayesian network 贝叶斯网络\quad 280\par
Bayesian probability 贝叶斯概率\quad 47\par
beam search 束搜索\quad 329\par
Bernoulli distribution 伯努利分布\quad 56, 81\par
Bernoulli mixture model 伯努利混合模型\quad 409\par
BERT BERT\quad 330\par
bi-gram model 二元模型\quad 323\par
bias 偏差\quad 35, 109\par
bias parameter 偏置参数\quad 98, 116, 156\par
bias-variance trade-off 偏差-方差权衡\quad 108\par
BigGAN BigGAN\quad 458\par
bijective function 双射函数\quad 466\par
binomial distribution 二项分布\quad 57\par
bits 比特\quad 41\par
blind source separation 盲源分离\quad 437\par
blocked path 阻塞路径\quad 290, 294\par
boosting 提升\quad 241\par
bootstrap 自助\quad 239\par
bottleneck 瓶颈问题\quad 325\par
bounding box 边界框\quad 265\par
Box-Muller method Box-Muller 方法\quad 368\par
byte pair encoding 字节对编码\quad 321\par
canonical correlation analysis 典型相关分析\quad 426\par
canonical link function 规范连接函数\quad 143\par
Cauchy distribution 柯西分布\quad 368\par
causal attention 因果注意力\quad 328\par
causality 因果\quad 297\par
central differences 中心差分法\quad 206\par
central limit theorem 中心极限定理\quad 60\par
ChatGPT ChatGPT\quad 335\par
child node 子节点\quad 216, 281\par
Cholesky decomposition Cholesky 分解\quad 368\par
circular normal distribution 圆形正态分布\quad 78\par
classical probability 经典概率\quad 47\par
classification 分类\quad 3\par
CLIP CLIP\quad 167\par
co-parents 共父节点\quad 298\par
codebook vector 码本向量\quad 338, 396\par
collider node 碰撞节点\quad 292\par
combining models 模型整合\quad 127\par
committee 委员会\quad 239\par
complete data set 完整数据集\quad 404\par
completing the square 完全平方\quad 66\par
computer vision 计算机视觉\quad 248\par
concave function 凹函数\quad 46\par
concentration parameter 聚焦参数\quad 78\par
condition number 条件数\quad 190\par
conditional entropy 条件熵\quad 47\par
conditional expectation 条件期望\quad 31\par
conditional independence 条件独立性\quad 127, 289\par
conditional mixture model 条件混合模型\quad 173\par
conditional probability 条件概率\quad 25\par
conditional VAE 条件 VAE\quad 491\par
conditioner 条件函数\quad 469\par
confusion matrix 混淆矩阵\quad 128\par
continuous bag of words 连续词袋\quad 320\par
continuous normalizing flow 连续标准化流\quad 475\par
contrastive divergence 对比散度\quad 387\par
contrastive learning 对比学习\quad 165\par
convex function 凸函数\quad 45\par
convolution 卷积\quad 249, 275\par
convolutional network 卷积网络\quad 247\par
correlation matrix 相关性矩阵\quad 428\par
cost function 成本函数\quad 124\par

% p.581
coupling flow 耦合流\quad 467\par
coupling function 耦合函数\quad 469\par
covariance 协方差\quad 32\par
Cox’s axioms Cox 定理\quad 48\par
cross attention 交叉注意力\quad 332\par
cross-correlation 互相关\quad 251, 275\par
cross-entropy error function 交叉熵误差函数\quad 139, 141, 170\par
cross-validation 交叉验证\quad 13\par
cumulative distribution function 累积分布函数\quad 29\par
curse of dimensionality 维度诅咒\quad 150\par
curve fitting 曲线拟合\quad 7\par
CycleGAN CycleGAN\quad 459\par
d-separation d 分离\quad 289, 293, 407\par
DAG, see directed acyclic graph DAG，参见“有向无环图”\par
data augmentation 数据增强\quad 166, 222\par
data compression 数据压缩\quad 396\par
DDIM DDIM\quad 505\par
DDPM DDPM\quad 493\par
decision 决策\quad 105\par
decision boundary 决策边界\quad 115, 122\par
decision region 决策区域\quad 115, 122\par
decision surface, see decision boundary 决策面，参见“决策边界”\par
decision theory 决策理论\quad 105, 121\par
decoder 解码器\quad 479\par
deep double descent 双重下降\quad 232\par
deep learning 深度学习\quad 17\par
deep neural networks 深度神经网络\quad 17\par
deep sets 深度集合\quad 356\par
DeepDream DeepDream\quad 264\par
degrees of freedom 自由度\quad 421\par
denoising 去噪\quad 493\par
denoising autoencoder 去噪自编码器\quad 482\par
denoising diffusion implicit model 去噪扩散隐式模型\quad 505\par
denoising diffusion probabilistic model 去噪扩散概率模型\quad 493\par
denoising score matching 去噪得分匹配\quad 507\par
density estimation 密度估计\quad 33, 55\par
dequantization 去量化\quad 448\par
descendant node 后代节点\quad 292\par
design matrix 设计矩阵\quad 101\par
development set 开发集\quad 13\par
diagonal covariance matrix 对角协方差矩阵\quad 64\par
differential entropy 微分熵\quad 44\par
diffusion kernel 扩散核\quad 495\par
diffusion model 扩散模型\quad 493\par
Dirac delta function 狄拉克 $\delta$ 函数\quad 30\par
directed acyclic graph 有向无环图\quad 282\par
directed cycle 有向环\quad 282\par
directed factorization 有向分解\quad 298\par
directed graph 有向图\quad 280\par
directed graphical model 有向图模型\quad 280\par
discriminant function 判别函数\quad 116, 125\par
discriminative model 判别模型\quad 126, 137, 297\par
disentangled representations 解纠缠表示\quad 462\par
distributed representation 分布式表示\quad 163\par
dot-product attention 点积注意力\quad 310\par
double descent 双重下降\quad 232\par
dropout Dropout\quad 241\par
E step E 步骤\quad 402, 405\par
early stopping 早停法\quad 230\par
earth mover’s distance 推土机距离\quad 457\par

% p.582
ECM, see expectation conditional maximization ECM，参见“期望条件最大化”\par
edge 边\quad 280, 349\par
edge detection 边缘检测\quad 252\par
ELBO, see evidence lower bound ELBO，参见“证据下界”\par
EM, see expectation maximization EM，参见“最大期望”\par
embedding space 嵌入空间\quad 163\par
embedding vector 嵌入向量\quad 349\par
encoder 编码器\quad 479\par
energy function 能量函数\quad 384\par
energy-based models 基于能量的模型\quad 384\par
ensemble methods 集成\quad 239\par
entropy 熵\quad 41\par
epistemic uncertainty 认知不确定性\quad 21\par
epoch 训练周期\quad 186\par
equality constraint 等式约束\quad 527\par
equivariance 等变性\quad 224, 251, 255, 317, 351\par
erf function erf 函数\quad 143\par
error backpropagation, see backpropagation 误差反向传播，参见“反向传播”\par
error function 误差函数\quad 8, 48, 169, 182\par
Euler-Lagrange equations 欧拉-拉格朗日方程\quad 525\par
evaluation trace 求值轨迹\quad 213\par
evidence lower bound 证据下界\quad 412, 439, 485, 499\par
expectation 期望\quad 31\par
expectation conditional maximization 期望条件最大化\quad 416\par
expectation maximization 最大期望\quad 400, 404, 440, 442\par
expectation step, see E step 期望步，参见“E 步骤”\par
expectations 期望\quad 366\par
explaining away 相消解释\quad 293\par
exploding gradient 梯度爆炸\quad 196, 325\par
exponential distribution 指数分布\quad 30, 367\par
exponential family 指数族分布\quad 80, 136, 282\par
expression swell 表达式膨胀\quad 212\par
factor analysis 因子分析\quad 436\par
factor graph 因子图\quad 280\par
factor loading 因子载荷\quad 436\par
false negative 假阴性\quad 23\par
false positive 假阳性\quad 23\par
fast gradient sign method 快速梯度符号\quad 263\par
fast R-CNN fast R-CNN\quad 270\par
feature extraction 特征提取\quad 17, 98\par
feature map 特征图\quad 251\par
features 特征\quad 155\par
feed-forward network 前馈网络\quad 149, 168\par
feed-forward networks 前馈神经网络\quad 16\par
few-shot learning 小样本学习\quad 165, 335\par
filter 滤波器\quad 250\par
fine-tuning 微调\quad 3, 19, 165, 334\par
flow matching 流匹配\quad 476\par
forward kinematics 正运动学\quad 172\par
forward problem 正问题\quad 172\par
forward propagation 前向传播\quad 203\par
foundation model 基础模型\quad 19, 306, 334, 349\par
frequentist probability 频率学派概率\quad 47\par
fuel system 燃油系统\quad 292\par
fully connected graphical model 全连接的图模型\quad 281\par
fully convolutional network 全卷积网络\quad 272\par

% p.583
functional 泛函\quad 524\par
Gabor filters Gabor 滤波器\quad 259\par
gamma distribution 伽马分布\quad 370\par
GAN, see generative adversarial network GAN，参见“生成对抗网络”\par
gated recurrent unit 门控循环单元\quad 326\par
Gaussian 高斯\quad 32, 59\par
Gaussian mixture 高斯混合\quad 74, 173, 234, 397\par
GEM, see generalized EM algorithm GEM，参见“广义 EM 算法”\par
generalization 泛化\quad 7\par
generalized EM algorithm 广义 EM 算法\quad 416\par
generalized linear model 广义线性模型\quad 138, 143\par
generative adversarial network 生成对抗网络\quad 453\par
generative AI 生成式 AI\quad 5\par
generative model 生成式模型\quad 5, 126, 296, 454\par
generative pre-trained transformer 生成式预训练 Transformer\quad 5, 326\par
geometric deep learning 几何深度学习\quad 362\par
Gibbs sampling 吉布斯采样\quad 380\par
global minimum 全局最小值\quad 183\par
GNN, see graph neural network GNN，参见“图神经网络”\par
GPT, see generative pre-trained transformer GPT，参见“生成式预训练 Transformer”\par
GPU, see graphics processing unit GPU，参见“图形处理单元”\par
gradient descent 梯度下降\quad 181\par
graph attention network 图注意力网络\quad 359\par
graph convolutional network 图卷积网络\quad 353\par
graph neural network 图神经网络\quad 347\par
graph representation learning 图表示学习\quad 349\par
graphical model 图模型\quad 279\par
graphical model factorization 图模型分解\quad 282\par
graphics processing unit 图形处理单元\quad 17, 306\par
group theory 群理论\quad 222\par
guidance 引导\quad 510\par
Hadamard product Hadamard 积\quad 468\par
Hamiltonian Monte Carlo 哈密顿蒙特卡洛\quad 384\par
handwritten digit 手写数字\quad 427\par
He initialization He 初始化\quad 188\par
head-to-head path 头头相接路径\quad 292\par
head-to-tail path 头尾相接路径\quad 291\par
Heaviside step function 单位阶跃函数\quad 140\par
Hessian matrix 黑塞矩阵\quad 183, 209\par
Hessian outer product approximation 黑塞外积近似法\quad 210\par
heteroscedastic 异方差\quad 173\par
hidden Markov model 隐马尔可夫模型\quad 324, 407\par
hidden unit 隐藏单元\quad 16, 156\par
hidden variable, see latent variable 隐变量，参见“潜变量”\par
hierarchical representation 层次化表示\quad 162\par
histogram density estimation 直方图密度估计\quad 85\par
history of machine learning 机器学习历史\quad 14\par
hold-out set 保留集\quad 13\par

% p.584
homogeneous Markov chain 均匀马尔可夫链\quad 377\par
Hooke’s law 胡克定律\quad 443\par
Hutchinson’s trace estimator Hutchinson 的迹估计器\quad 474\par
hybrid Monte Carlo 混合蒙特卡洛\quad 384\par
hyperparameter 超参数\quad 12\par
IAF, see inverse autoregressive flow IAF，参见“逆自回归流”\par
ICA, see independent component analysis ICA，参见“独立成分分析”\par
identifiability 可识别性\quad 400\par
IID, see independent and identically distributed IID，参见“独立同分布”\par
image segmentation 图像分割\quad 270\par
ImageNet data set ImageNet 数据集\quad 257\par
importance sampling 重要性采样\quad 371, 383\par
importance weight 重要性权重\quad 372\par
improper distribution 反常分布\quad 30\par
improper prior 反常先验\quad 227\par
inactive constraint 非激活的约束\quad 528\par
incomplete data set 不完整的数据集\quad 404\par
independent and identically distributed 独立同分布\quad 37, 294\par
independent component analysis 独立成分分析\quad 437\par
independent factor analysis 独立因子分析\quad 438\par
independent variables 独立变量\quad 28\par
inductive bias 归纳偏置\quad 17, 220\par
inductive learning 归纳学习\quad 349, 358\par
inequality constraint 不等式约束\quad 527\par
inference 推理\quad 105, 121, 125, 288\par
InfoNCE InfoNCE\quad 166\par
information theory 信息论\quad 46\par
instance discrimination 个体判别\quad 166\par
internal covariate shift 内部协变量偏移\quad 197\par
internal representation 内部表示\quad 265\par
intersection-over-union 交并比\quad 266\par
intrinsic dimensionality 内在维度\quad 421\par
invariance 不变性\quad 222, 255, 351\par
inverse autoregressive flow 逆自回归流\quad 471\par
inverse kinematics 逆运动学\quad 172\par
inverse problem 逆问题\quad 107, 172, 220, 296\par
Iris data 鸢尾花数据\quad 151\par
IRLS, see iterative reweighted least squares IRLS，参见“迭代重加权最小二乘法”\par
isotropic covariance matrix 各向同性协方差矩阵\quad 64\par
iterative reweighted least squares 迭代重加权最小二乘法\quad 140\par
Jacobian matrix 雅可比矩阵\quad 39, 207\par
Jensen’s inequality 詹森不等式\quad 46\par
Jensen-Shannon divergence Jensen-Shannon 散度\quad 463\par
K nearest neighbours K 近邻\quad 89\par
K-means clustering algorithm K 均值聚类算法\quad 392, 408\par
Kalman filter 卡尔曼滤波器\quad 302, 439\par
Karush-Kuhn-Tucker conditions 卡鲁什-库恩-塔克条件\quad 528\par
kernel density estimator 核密度估计\quad 86, 506\par
kernel function 核函数\quad 87\par
kernel image 核图像\quad 250\par
KKT, see Karush-Kuhn-Tucker conditions KKT，参见“卡鲁什-库恩-塔克条件”\par

% p.585
KL divergence, see Kullback-Leibler divergence KL 散度，参见“Kullback-Leibler 散度”\par
Kosambi-Karhunen-Loève transform Kosambi-Karhunen-Loève 变换\quad 422\par
Kullback-Leibler divergence Kullback-Leibler 散度\quad 45, 463\par
Lagrange multiplier 拉格朗日乘子\quad 526\par
Lagrangian 拉格朗日函数\quad 527\par
Langevin dynamics 朗之万动力学\quad 386\par
Langevin sampling 朗之万采样\quad 387\par
language model 语言模型\quad 326\par
Laplace distribution 拉普拉斯分布\quad 30\par
large language model 大语言模型\quad 5, 326, 333\par
lasso lasso\quad 229\par
latent class analysis 潜在类别分析\quad 409\par
latent diffusion model 潜扩散模型\quad 512\par
latent variable 潜变量\quad 64, 288, 391, 421\par
layer normalization 层归一化\quad 198, 315\par
LDM, see latent diffusion model LDM，参见“潜扩散模型”\par
LDS, see linear dynamical system LDS，参见“线性动态系统”\par
leaky ReLU leaky ReLU\quad 161\par
learning curve 学习曲线\quad 192, 230\par
learning rate parameter 学习率\quad 192\par
learning to learn 学会学习\quad 165\par
least-mean-squares algorithm 最小二乘\quad 103\par
least-squares GAN 最小二乘 GAN\quad 456\par
leave-one-out 留一法\quad 13\par
LeNet convolutional network LeNet\quad 257\par
Levenberg-Marquardt approximation 列文伯格-马夸尔特近似法\quad 210\par
likelihood function 似然函数\quad 34, 399\par
likelihood weighted sampling 似然加权采样\quad 383\par
linear discriminant 线性判别式\quad 116\par
linear dynamical system 线性动态系统\quad 439\par
linear independence 线性无关\quad 518\par
linear regression 线性回归\quad 6, 97\par
linear-Gaussian model 线性高斯模型\quad 67, 284\par
linearly separable 线性可分\quad 115\par
link, see edge 链接，参见“边”\par
link function 连接函数\quad 137, 144\par
LLM, see large language model LLM，参见“大语言模型”\par
LMS, see least-mean-squares algorithm LMS，参见“平均最小均方算法”\par
local minimum 局部极小值\quad 182\par
log odds 对数几率\quad 132\par
logic sampling 逻辑采样\quad 383\par
logistic regression 逻辑斯谛回归\quad 138\par
logistic sigmoid 逻辑斯谛 sigmoid\quad 81, 99, 131, 139\par
logit function logit 函数\quad 132\par
long short-term memory 长短时记忆\quad 326\par
LoRA, see low-rank adaptation LoRA，参见“低秩自适应”\par
loss function 损失函数\quad 106, 124\par
loss matrix 损失矩阵\quad 124\par
lossless data compression 无损数据压缩\quad 396\par
lossy data compression 有损数据压缩\quad 396\par
low-rank adaptation 低秩自适应\quad 335\par
LSGAN, see least-squares GAN LSGAN，参见“最小二乘 GAN”\par
LSTM, see long short-term memory LSTM，参见“长短时记忆”\par
